\documentclass[12pt,x11names,a4paper]{article}
\input{preamble}


\newgeometry{margin=2cm}

\pagestyle{fancy}
\fancyhf{}

\rhead{Nørre Gymnasium\\1.b
}
\cfoot{Side \thepage \hspace{1pt} af \pageref{LastPage}}

%Husk at rette modul og dato!
\lhead{Delprøve 1 - prøveeksamen\\ Matematik B
}
\chead{29. februar
}

\begin{document}

%\includepdf[pages=-]{Forsider/aarsprove_1v.pdf}
\savegeometry{art}

\begin{titlepage}
\newgeometry{margin=0pt}

\begin{minipage}{0.27\textwidth}

\begin{tikzpicture}[overlay]
\fill[top color = NorregGroen!40, bottom color = NorregGroen] (6,10) rectangle (-10,-30);
\end{tikzpicture}
\end{minipage}
\begin{minipage}{0.73\textwidth}
\begin{center}
\phantom{h} \vspace{1cm}\\
\hspace{4cm}
\includegraphics[scale = 1]{Billeder/Norreg.png} \\
\phantom{h} \vspace{5cm}\\
\rule{0.7\textwidth}{0.3mm}\\
\phantom{h}\\
{\fontsize{50}{60}\selectfont Delprøve 1 \\ Prøveeksamen}\\
\phantom{h}\\
\rule{0.7\textwidth}{0.3mm}\\
\Large 2024\\
\Large 1.b Ma

\end{center}
\end{minipage}
\end{titlepage}
\loadgeometry{art}

%Udfyld afsnit herunder og lav til egen Latex-fil

%Kopier følgende til overskrift:

%\begin{center}
%\Huge
%Aflevering 1
%\end{center}
%\section*{Opgave 1}
%\stepcounter{section}
\begin{center}
%Opgavesætter er delt i to dele:\\
%Delprøve 1 kun med den centralt udmeldte formelsamling.\\
%Delprøve 2 med alle hjælpemidler.
\end{center}

\begin{comment}
\section*{Krav til formidling af din besvarelse}

Ved bedømmelse af helhedsindtrykket af besvarelsen af de enkelte opgaver lægges særlig vægt på følgende fire punkter:
\begin{itemize}
\item[$\cdot$] \textbf{Redegørelse og dokumentation for metode} \\
Besvarelsen skal indeholde en redegørelse for den anvendte løsningsstragegi med dokumentation i form af et passende antal mellemregninger \textit{eller} matematiske forklaringer på metoden, når et matematisk værktøjsprogram anvendes.
\item[$\cdot$] \textbf{Figurer, grafer og andre illustrationer} \\
Besvarelsen skal indeholde hensigtsmæssig brug af figurer, grafer og andre illustrationer, og der skal være tydelige henvisninger til brug af disse i den forklarende tekst.
\item[$\cdot$] \textbf{Notation og layout}\\
Besvarelsen skal i overensstemmelse med god matematisk skik opstilles med hensigtsmæssig brug af symbolsprog, og med en redegørelse for den matematiske notation, der indføres og anvendes, og som ikke kan henføres stil standardviden.
\item[$\cdot$] \textbf{Formidling og forklaring}\\
Besvarelsen af rene matematikopgaver skal indeholde en angivelse af givne oplysninger og korte forklaringer knyttet til den anvendte løsningsstrategi beskrevet med brug af almindelig matematisk notation. 

Besvarelsen af opgaver, der omhandler matematiske modeller, skal indeholde en kort præsentation af modellens kontekst, herunder betydning af modellens parametre. De enkelte delspørgsmål skal afsluttes med en præcis konklusion præsenteret i et klart sprog i relation til konteksten.
\end{itemize}

\newpage
\end{comment}


\begin{opgavetekst}{Opgave 1}
	\begin{center}
		\includegraphics[width = 0.6\textwidth]{Billeder/Zoo1}
	\end{center}
	Den zoologiske have Metropol Zoo har hvert år siden år 1980 målt antallet af besøgende. Deres målinger 
	fremgår af \href{https://github.com/ChristianJLex/TeachingNotes/raw/master/2023-2024/Data og lign/BesoegendeiZoo.xlsx}{\color{blue!60} dette datasæt}.
	Metropol Zoo har en forventning om, at antallet af besøgende $f$ (i tusinde) kan
	beskrives ved forskriften
	\begin{align*}
		f(x) = b\cdot x^a,
	\end{align*}
	hvor $x$ betegner antal år efter år 1980.
\end{opgavetekst}


\begin{delopgave}{}{1}
	Brug tallene fra datasættet til at bestemme tallene $a$ og $b$, og opskriv forskriften for 
	modellen, der beskriver antallet af besøgende.
\end{delopgave}
\begin{delopgave}{}{2}
	Bestem antallet af besøgende i år 2030 i følge modellen. 
\end{delopgave}
\begin{meretekst}
	Da initiativtagerne bag Metropol Zoo opstartede zooen, så havde de en forventning
	om, at antallet af besøgende kunne beskrives ved sammenhængen
	\begin{align*}
		g(x) = 4x + 200,
	\end{align*}
	hvor $g$ betegner det årlige antal besøgende i tusinde og $x$ betegner antal år efter år 
	1980
\end{meretekst}
\begin{delopgave}{}{3}
	Tegn forskriften for $f$ og $g$ i samme koordinatsystem på intervallet $[1,100]$. 
\end{delopgave}
\begin{delopgave}{}{4}
	Bestem det årstal, hvor $g$ forudser, at der vil være flere besøgende i Metropol Zoo end
	modellen $f$.
\end{delopgave}

%%%%%%%%%%%%%%%%%%%%%%%%%%%%%%%%%%%%%%%%%%%%%%%%%%%%%%%%%%%%%%%%%%%%%%%%%%%%%%%%%%%%%%%%%%%%%%%%
%%%%%%%%%%%%%%%%%%%%%%%%%%%%%%%%%%%%%%%%%%%%%%%%%%%%%%%%%%%%%%%%%%%%%%%%%%%%%%%%%%%%%%%%%%%%%%%%
%%%%%%%%%%%%%%%%%%%%%%%%%%%%%%%%%%%%%%%%%%%%%%%%%%%%%%%%%%%%%%%%%%%%%%%%%%%%%%%%%%%%%%%%%%%%%%%%
%%%%%%%%%%%%%%%%%%%%%%%%%%%%%%%%%%%%%%%%%%%%%%%%%%%%%%%%%%%%%%%%%%%%%%%%%%%%%%%%%%%%%%%%%%%%%%%%
%%%%%%%%%%%%%%%%%%%%%%%%%%%%%%%%%%%%%%%%%%%%%%%%%%%%%%%%%%%%%%%%%%%%%%%%%%%%%%%%%%%%%%%%%%%%%%%%
%%%%%%%%%%%%%%%%%%%%%%%%%%%%%%%%%%%%%%%%%%%%%%%%%%%%%%%%%%%%%%%%%%%%%%%%%%%%%%%%%%%%%%%%%%%%%%%%
%%%%%%%%%%%%%%%%%%%%%%%%%%%%%%%%%%%%%%%%%%%%%%%%%%%%%%%%%%%%%%%%%%%%%%%%%%%%%%%%%%%%%%%%%%%%%%%%
%%%%%%%%%%%%%%%%%%%%%%%%%%%%%%%%%%%%%%%%%%%%%%%%%%%%%%%%%%%%%%%%%%%%%%%%%%%%%%%%%%%%%%%%%%%%%%%%
\begin{opgavetekst}{Opgave 2}
	\begin{center}
		\includegraphics[width=0.6\textwidth]{Billeder/penge}
	\end{center}
	Metropol Zoo ønsker sig et bæverhabitat i haven. De har ikke råd til 
	at betale for habitatet, så de skal have et lån i banken. Banken fortæller 
	dem at beløbet de skylder banken vil kunne beskrives ved funktionen 
	\begin{align*}
		h(x) = 19.2\cdot 1.052^x,
	\end{align*}
	hvis de ikke betaler af på deres lån. $h$ betegner det skyldige beløb i mio. kr.,
	og $x$ betegner antal år efter optagelsen af lånet. 
\end{opgavetekst}
\begin{delopgave}{}{1}
	Forklar, hvad tallene 19.2 og 1.052 siger om banklånet.
\end{delopgave}
\begin{delopgave}{}{2}
	Afgør, hvor længe der vil gå før Metropol Zoo skylder 30 mio., hvis de ikke betaler af på
	deres lån.
\end{delopgave}
\begin{meretekst}
	Metropol Zoo skal selvfølgelig betale af på deres lån. Deres lån kan beskrives ved formlen
	\begin{align*}
		19.2 = y\cdot \frac{1 - (1+r)^{-n}}{r}.
	\end{align*}
	For formlen gælder der, at 
	\begin{itemize}
		\item[$\cdot$] $y$ er beløbet, der afdrages årligt på lånet (i mio. kr.).
		\item[$\cdot$] $r$ er renten på lånet. 
		\item[$\cdot$] $n$ er antal terminer på lånet - altså i dette tilfælde antal år, som 
		lånet skal løbe i.  	
	\end{itemize}
	\phantom{h}
\end{meretekst}	
\begin{delopgave}{}{3}
	Indsæt renten fra a) og bestem, hvad Metropol Zoo skal betale tilbage hvert år, hvis de 
	skal betale lånet af på 30 år.
\end{delopgave}
\begin{delopgave}{}{4}
	Afgør hvornår lånet er betalt tilbage, hvis de afdrager 1 mio. kr. om året på lånet. 
\end{delopgave}

%%%%%%%%%%%%%%%%%%%%%%%%%%%%%%%%%%%%%%%%%%%%%%%%%%%%%%%%%%%%%%%%%%%%%%%%%%%%%%%%%%%%%%%%%%%%%%%%
%%%%%%%%%%%%%%%%%%%%%%%%%%%%%%%%%%%%%%%%%%%%%%%%%%%%%%%%%%%%%%%%%%%%%%%%%%%%%%%%%%%%%%%%%%%%%%%%
%%%%%%%%%%%%%%%%%%%%%%%%%%%%%%%%%%%%%%%%%%%%%%%%%%%%%%%%%%%%%%%%%%%%%%%%%%%%%%%%%%%%%%%%%%%%%%%%
%%%%%%%%%%%%%%%%%%%%%%%%%%%%%%%%%%%%%%%%%%%%%%%%%%%%%%%%%%%%%%%%%%%%%%%%%%%%%%%%%%%%%%%%%%%%%%%%
%%%%%%%%%%%%%%%%%%%%%%%%%%%%%%%%%%%%%%%%%%%%%%%%%%%%%%%%%%%%%%%%%%%%%%%%%%%%%%%%%%%%%%%%%%%%%%%%
%%%%%%%%%%%%%%%%%%%%%%%%%%%%%%%%%%%%%%%%%%%%%%%%%%%%%%%%%%%%%%%%%%%%%%%%%%%%%%%%%%%%%%%%%%%%%%%%
%%%%%%%%%%%%%%%%%%%%%%%%%%%%%%%%%%%%%%%%%%%%%%%%%%%%%%%%%%%%%%%%%%%%%%%%%%%%%%%%%%%%%%%%%%%%%%%%
%%%%%%%%%%%%%%%%%%%%%%%%%%%%%%%%%%%%%%%%%%%%%%%%%%%%%%%%%%%%%%%%%%%%%%%%%%%%%%%%%%%%%%%%%%%%%%%%
\newpage
\begin{opgavetekst}{Opgave 3}
	\begin{center}
		\includegraphics[width=0.6\textwidth]{Billeder/beaver}
	\end{center}
	En fuldvoksen europæisk bæver spiser omkring 1.8kg plantemateriale om dagen, der 
	hovedsagligt består af græsser, vandplanter og bark. Metropol Zoo vil gerne have 8 bævere 
	i haven.
\end{opgavetekst}
\begin{delopgave}{}{1}
	Opstil en model, der beskriver mængden af foder til bæverne som Metropol Zoo har forbrugt
	som funktion af antal forløbne dage. 
\end{delopgave}
\begin{delopgave}{}{2}
	Brug modellen til at bestemme hvor meget bæverfoder, haven forbruger på et år.
\end{delopgave}
\begin{meretekst}
	For både føde og materiale til dæmninger fælder den europæiske bæver træer på op til 60cm i
	diameter. Tiden, det tager for bæveren at fælde et aspetræ er givet ved
	\begin{align*}
		f(x) = 0.00162\cdot \pi \cdot x^3,
	\end{align*}
	hvor $f$ er tiden i timer og $x$ er radius på stammen i cm. 
\end{meretekst}
\begin{delopgave}{}{3}
	Afgør, hvilken funktionstype $f$ er, og bestem tallene $a$ og $b$. 
\end{delopgave}
\begin{delopgave}{}{4}
	Bestem hvor mange gange længere det vil tage for en bæver at fælde et træ, hvis radius af 
	træet fordobles.
\end{delopgave}
%%%%%%%%%%%%%%%%%%%%%%%%%%%%%%%%%%%%%%%%%%%%%%%%%%%%%%%%%%%%%%%%%%%%%%%%%%%%%%%%%%%%%%%%%%%%%%%%
%%%%%%%%%%%%%%%%%%%%%%%%%%%%%%%%%%%%%%%%%%%%%%%%%%%%%%%%%%%%%%%%%%%%%%%%%%%%%%%%%%%%%%%%%%%%%%%%
%%%%%%%%%%%%%%%%%%%%%%%%%%%%%%%%%%%%%%%%%%%%%%%%%%%%%%%%%%%%%%%%%%%%%%%%%%%%%%%%%%%%%%%%%%%%%%%%
%%%%%%%%%%%%%%%%%%%%%%%%%%%%%%%%%%%%%%%%%%%%%%%%%%%%%%%%%%%%%%%%%%%%%%%%%%%%%%%%%%%%%%%%%%%%%%%%
%%%%%%%%%%%%%%%%%%%%%%%%%%%%%%%%%%%%%%%%%%%%%%%%%%%%%%%%%%%%%%%%%%%%%%%%%%%%%%%%%%%%%%%%%%%%%%%%
%%%%%%%%%%%%%%%%%%%%%%%%%%%%%%%%%%%%%%%%%%%%%%%%%%%%%%%%%%%%%%%%%%%%%%%%%%%%%%%%%%%%%%%%%%%%%%%%
%%%%%%%%%%%%%%%%%%%%%%%%%%%%%%%%%%%%%%%%%%%%%%%%%%%%%%%%%%%%%%%%%%%%%%%%%%%%%%%%%%%%%%%%%%%%%%%%
%%%%%%%%%%%%%%%%%%%%%%%%%%%%%%%%%%%%%%%%%%%%%%%%%%%%%%%%%%%%%%%%%%%%%%%%%%%%%%%%%%%%%%%%%%%%%%%%
\newpage
\begin{opgavetekst}{Opgave 4}
	\begin{center}
		\includegraphics[width=0.6\textwidth]{Billeder/tigerunge}
	\end{center}
	En løvinde hos Metropol Zoo har født løveungen Jørgen. De håber på, at dette kan øge
	 deres omsætning med 15 procent det næste år. De har dog kun en forventning
	 om, at løveungen vil øge deres besøgstal med 8 procent. 
\end{opgavetekst}
\begin{delopgave}{}{1}
	Bestem hvor meget Metropol Zoo skal øge deres billetpriser med, hvis de skal opnå en
	stigning i omsætningen på 15$\%$. 
\end{delopgave}
\begin{meretekst}
	Vægten af løveungen beskrives ved funktionen
	\begin{align*}
		V(t) = 260\cdot(1-e^{-0.082(t+3.54)}),
	\end{align*}
	hvor $V$ er vægten i kg og $t$ er løvens alder i måneder. 
\end{meretekst}
%
\begin{delopgave}{}{2}
	Tegn grafen for $V$ på et passende interval, og se om du ud fra grafen kan
	gennemskue hvad tallet 260 betyder for modellen. 
\end{delopgave}
%
\begin{delopgave}{}{3}
	Bestem hvad Jørgen vejer efter ét år i følge modellen.
\end{delopgave}

\end{document}




